\section{问题重述}

\subsection{问题背景}

碳化硅是一种综合性能优越的第三代半导体材料，其外延层的厚度是外延材料的关键参数之一，对器件性能有重要影响。制定一套科学、准确、可靠的外延层厚度测试方法，具有重要的工程意义。


红外干涉法是外延层厚度测量的无损伤测量方法。外延层与衬底因掺杂载流子浓度的不同而有不同的折射率，红外光入射后，一部分从外延层表面反射出来，另一部分在外延层与衬底的界面反射回来，两束光在一定条件下产生干涉条纹。根据红外光谱的波长、外延层的折射率和红外光的入射角等参数，即可确定外延层的厚度。


通常外延层的折射率不是常数，它与掺杂载流子的浓度、红外光谱的波长等参数有关。光波在两界面间还可能发生多次反射与透射，产生多光束干涉，进而影响厚度的测定。本文围绕上述三个层次，依次建立两光束干涉模型、设计厚度求解算法，并研究多光束干涉的判别与消除。

\subsection{需要解决的问题}

本文需要解决以下三个问题。


{\textbf{问题一}}：考虑外延层和衬底界面只有一次反射、透射所产生干涉条纹的情形，建立确定外延层厚度的数学模型。


{\textbf{问题二}}：根据问题一的数学模型，设计确定外延层厚度的算法。对附件 1 和附件 2 提供的碳化硅晶圆片实测光谱（入射角分别为 $10^\circ$ 和 $15^\circ$，第 1 列为波数，第 2 列为反射率）给出计算结果，并分析结果的可靠性。


{\textbf{问题三}}：光波在外延层界面和衬底界面可以产生多次反射和透射，从而产生多光束干涉。推导产生多光束干涉的必要条件，分析多光束干涉对外延层厚度计算精度可能产生的影响。根据必要条件，分析附件 3 和附件 4 提供的硅晶圆片测试结果（入射角分别为 $10^\circ$ 和 $15^\circ$）是否出现多光束干涉，给出确定硅外延层厚度的数学模型、算法及相应的计算结果；若认为多光束干涉也会出现在碳化硅晶圆片的测试结果中并影响厚度计算的精度，则设法消除其影响，并给出消除影响后的计算结果。
